\documentclass[11pt]{article}
\usepackage[T1]{fontenc}
\usepackage{lmodern}
\usepackage{amsmath,amssymb,amsthm,mathtools}
\usepackage[a4paper,margin=1in]{geometry}
\usepackage[colorlinks=true,linkcolor=blue,urlcolor=blue]{hyperref}
\newtheorem{theorem}{Theorem}[section]
\newtheorem{lemma}[theorem]{Lemma}
\newtheorem{corollary}[theorem]{Corollary}
\newtheorem{remark}[theorem]{Remark}
\newcommand{\Td}{\mathbb T^2}
\newcommand{\dd}{\,dx}
\newcommand{\norm}[1]{\left\|#1\right\|}
\title{\textbf{Fixed-Mass Nodal Coercivity for Density-Weighted Curvature}\\[4pt]
\large An Independent Auxiliary Verification Note for WCT Computation}
\author{Richard J. Reyes}
\date{October 27, 2026}
\begin{document}
\maketitle
\begin{abstract}
On a fixed two-dimensional periodic domain, we prove that positive conserved mass
and an upper bound on the cubic--quintic projection-free Hamiltonian enforce a
strictly positive and quantitatively explicit lower comparison between
density-weighted Laplacian curvature and ordinary Laplacian energy.
No uniform positive lower bound on pointwise amplitude is imposed, so fields
with nodes are allowed. This closes a previously open \emph{fixed-domain,
fixed-resource} analytic step in a separate curvature--communication study.
The constant is not controlled uniformly in domain size, spectral resolution,
or computational input length. An exact two-site version and the algebraic
readout-cost transfer are formalized in Lean. This auxiliary note is independent
of the author's original P versus NP manuscript, which it neither revises nor
claims to complete.
\end{abstract}

\section{The maintained Hamiltonian and the curvature weight}
Let \(\Omega=\mathbb R^2/\Lambda\) be a fixed connected flat two-dimensional
torus, of volume \(V=|\Omega|\), and let \(u\in H^2(\Omega;\mathbb C)\).
For \(\kappa=k_\star^2\geq0\) and \(\beta,\gamma>0\), define
\begin{equation}
\label{eq:H}
\mathcal H[u]=\int_\Omega
\left(|(\Delta+\kappa)u|^2-\frac{\beta}{2}|u|^4
+\frac{\gamma}{3}|u|^6\right)\dd.
\end{equation}
This is the Hamiltonian of the separately maintained projection-free evolution
\begin{equation}
\label{eq:pde}
i\partial_tu=(\Delta+\kappa)^2u-\beta|u|^2u+\gamma|u|^4u.
\end{equation}
The present theorem concerns the \emph{distinct} density-weighted curvature
functional, not the Hamiltonian's own quadratic term. For
\(\varepsilon>0\), \(\alpha\geq0\), let
\begin{equation}
\label{eq:Q}
F_{\varepsilon,\alpha}(s)=
\frac{s}{(s+\varepsilon^2e^{-2\alpha s})^2},\qquad
\mathcal Q_{\varepsilon,\alpha}[u]
=\int_\Omega F_{\varepsilon,\alpha}(|u|^2)|\Delta u|^2\dd.
\end{equation}
The weight vanishes at every node \(u=0\). Pointwise coercivity is therefore
false without a node-free assumption, and the result below is genuinely global.

\section{Conservation-derived compactness}
Put \(a_{\beta,\gamma}=3\beta^2/(16\gamma)\). For every \(s\geq0\),
\begin{equation}
\label{eq:potential}
-\frac{\beta}{2}s^2+\frac{\gamma}{3}s^3
+a_{\beta,\gamma}s
=\frac{\gamma}{3}s
\left(s-\frac{3\beta}{4\gamma}\right)^2\geq0.
\end{equation}
Consequently
\begin{equation}
\label{eq:band}
\norm{(\Delta+\kappa)u}_{2}^{2}
\leq\mathcal H[u]+a_{\beta,\gamma}\norm{u}_{2}^{2}.
\end{equation}
Given a positive mass \(m>0\) and finite \(H_\star\in\mathbb R\), define
\[
\mathcal A_{m,H_\star}=
\{u\in H^2(\Omega;\mathbb C):
\norm{u}_2^2=m,\quad\mathcal H[u]\leq H_\star\}.
\]
If this class is nonempty, \(H_\star+a_{\beta,\gamma}m\geq0\).
For every \(u\in\mathcal A_{m,H_\star}\),
\begin{equation}
\label{eq:lapbound}
\norm{\Delta u}_2
\leq\sqrt{H_\star+a_{\beta,\gamma}m}+\kappa\sqrt m=:B.
\end{equation}
On the fixed torus, elliptic regularity gives
\(\norm{u}_{H^2}\leq C_\Omega(\sqrt m+B)\).
The compact embedding \(H^2(\Td)\Subset C^0(\Td)\)
therefore gives a common amplitude bound \(\norm{u}_\infty\leq U\)
and subsequential uniform convergence.

\section{The fixed-mass nodal theorem}
\begin{theorem}[Uniform nodal coercivity on a fixed Hamiltonian sublevel]
\label{thm:coercivity}
Fix \(\Omega,\kappa,\beta,\gamma,m,H_\star,\varepsilon\) as above, with
\(\mathcal A_{m,H_\star}\ne\varnothing\). There exists a number
\(c_\star=c_\star(\Omega,\kappa,\beta,\gamma,m,H_\star,\varepsilon)>0\)
such that for \emph{every} \(\alpha\geq0\) and every
\(u\in\mathcal A_{m,H_\star}\),
\begin{equation}
\label{eq:coercive}
\boxed{\qquad
\mathcal Q_{\varepsilon,\alpha}[u]
\geq c_\star\norm{\Delta u}_2^2.\qquad}
\end{equation}
No hypothesis \(\inf_\Omega|u|>0\) is needed.
A constructive coefficient is supplied in
Theorem~\ref{thm:explicit}; it depends on the fixed domain spectrum
and the resource upper bound.
\end{theorem}
\begin{proof}
From \(\alpha\geq0\) and the common amplitude bound,
\[
0\leq|u|^2+\varepsilon^2e^{-2\alpha|u|^2}
\leq U^2+\varepsilon^2,
\]
and therefore
\begin{equation}
\label{eq:product}
\mathcal Q_{\varepsilon,\alpha}[u]\geq
\frac{\norm{u\,\Delta u}_2^2}{(U^2+\varepsilon^2)^2}.
\end{equation}
Suppose no positive \(c_\star\) works. Then for each integer \(j\geq1\)
there are \(u_j\in\mathcal A_{m,H_\star}\) and \(\alpha_j\geq0\)
with \(\norm{\Delta u_j}_2>0\) and
\[
\mathcal Q_{\varepsilon,\alpha_j}[u_j]
<j^{-1}\norm{\Delta u_j}_2^2.
\]
The uniform bound \eqref{eq:lapbound} implies that the left side tends to
zero. By \eqref{eq:product},
\(\norm{u_j\Delta u_j}_2\to0\). Periodic integration by parts then yields
\[
\norm{\nabla u_j}_2^2
=-\operatorname{Re}\int_\Omega\overline{u_j}\Delta u_j\dd
\leq V^{1/2}\norm{u_j\Delta u_j}_2\longrightarrow0.
\]
By the Poincar\'e inequality,
\(u_j-\langle u_j\rangle\to0\) in \(L^2\), where
\(\langle u\rangle=V^{-1}\int_\Omega u\dd\).
Pass to a subsequence that converges uniformly in \(C^0\), using the
\(H^2\) compactness established above. Its limit is spatially constant,
and fixed mass forces the modulus of that constant to equal
\(\sqrt{m/V}>0\).

Thus, for all sufficiently large \(j\), the \emph{whole} field satisfies
\(|u_j(x)|\geq\frac12\sqrt{m/V}\), for every \(x\in\Omega\).
Together with the common upper amplitude bound,
\[
F_{\varepsilon,\alpha_j}(|u_j|^2)
\geq\frac{m}{4V(U^2+\varepsilon^2)^2}
=:c_0>0
\qquad\text{at every }x\in\Omega.
\]
Hence
\(\mathcal Q_{\varepsilon,\alpha_j}[u_j]
\geq c_0\norm{\Delta u_j}_2^2\),
contradicting its defining strict inequality for all \(j>c_0^{-1}\).
\end{proof}

\begin{remark}[What fixed mass does and does not imply]
The theorem does \emph{not} prove that each fixed-mass field is node-free.
It proves a uniform \emph{integrated} comparison over a compactly controlled
class of fields. The positive mass is essential: the earlier high-frequency
sequence \(N^{-2}\sin(Nx)\) has mass tending to zero and thus does not
contradict Theorem~\ref{thm:coercivity}. The Hamiltonian upper bound is
essential to prevent the Laplacian energy of arbitrarily fine oscillations
from escaping the compact class. The same compactness argument works
in dimensions \(d\leq3\) on a fixed flat torus; \(d=4\) requires
different control because \(H^2\not\hookrightarrow C^0\) there.
\end{remark}


\section{An explicit spectral coercivity coefficient}
\label{sec:explicit}
The compactness argument above can be made quantitative on the
fixed torus. Let \(\Lambda^\ast\) be the reciprocal lattice, and use
the orthonormal Fourier basis \(V^{-1/2}e^{ik\cdot x}\).
Let
\[
B_+=1+\sqrt{H_\star+a_{\beta,\gamma}m}+\kappa\sqrt m>0,\qquad
\mathcal S_\Omega=\frac1{\sqrt V}
\left(\sum_{k\in\Lambda^\ast}\frac{1}{1+|k|^4}\right)^{1/2},
\quad
U_+=\mathcal S_\Omega\sqrt{m+B_+^2}.
\]
Absolute convergence holds because the torus dimension is two.
Parseval and Cauchy--Schwarz show, without invoking an unspecified
embedding constant, that every \(u\in\mathcal A_{m,H_\star}\)
satisfies
\(\norm{\Delta u}_2\le B_+\) and \(\norm{u}_\infty\le U_+\).

Let \(\lambda_1>0\) be the first nonzero eigenvalue of
\(-\Delta\), and set \(\tau=\sqrt{m/(32V)}\).
Choose a finite radius \(R\) containing at least one nonzero reciprocal
lattice point, sufficiently large that
\[
\mathsf H_R B_+\le\tau,\qquad
\mathsf H_R=\frac1{\sqrt V}
\left(\sum_{|k|>R}|k|^{-4}\right)^{1/2}.
\]
Such an \(R\) exists since the Fourier tail converges.
Define
\[
\mathsf L_R=\frac1{\sqrt V}
\left(\sum_{0<|k|\le R}|k|^{-2}\right)^{1/2}>0,
\quad
\eta=\min\left\{\sqrt{\frac{\lambda_1 m}{2}},
\frac{\tau}{\mathsf L_R}\right\}>0.
\]

\begin{theorem}[Explicit fixed-domain coercivity coefficient]
\label{thm:explicit}
The constant in Theorem~\ref{thm:coercivity} can be chosen to be
\begin{equation}
\label{eq:explicit-c}
\boxed{\quad
c_{\mathrm{exp}}=
\frac{1}{(U_+^2+\varepsilon^2)^2}
\min\left\{\frac{m}{8V},
\frac{\eta^4}{VB_+^2}\right\}>0.\quad}
\end{equation}
In particular this choice is uniform for \emph{all}
\(\alpha\ge0\), and makes the dependence on the domain's Fourier
spectrum, positive mass, and energy upper bound explicit.
\end{theorem}
\begin{proof}
Write \(g=\norm{\nabla u}_2\) and \(\bar u=V^{-1}\int_\Omega u\dd\).
We distinguish two cases.

If \(g\le\eta\), Poincar\'e and the mass identity give
\[
V|\bar u|^2=m-\norm{u-\bar u}_2^2
\ge m-\lambda_1^{-1}g^2\ge m/2.
\]
Splitting the Fourier series of \(u-\bar u\) at radius \(R\)
and applying Cauchy--Schwarz to the gradient and Laplacian terms gives
\[
\norm{u-\bar u}_\infty
\le \mathsf L_R g+\mathsf H_R\norm{\Delta u}_2
\le2\tau=\sqrt{\frac{m}{8V}}.
\]
Since \(|\bar u|\ge\sqrt{m/(2V)}\),
the reverse triangle inequality yields
\(|u(x)|\ge\sqrt{m/(8V)}\) for every \(x\).
Consequently
\[
\mathcal Q_{\varepsilon,\alpha}[u]
\ge\frac{m}{8V(U_+^2+\varepsilon^2)^2}
     \norm{\Delta u}_2^2.
\]

If \(g>\eta\), periodic integration by parts, Cauchy--Schwarz,
and the uniform denominator estimate instead yield
\[
\norm{u\,\Delta u}_2^2\ge\frac{g^4}{V}
>\frac{\eta^4}{V},
\qquad
\mathcal Q_{\varepsilon,\alpha}[u]
\ge\frac{\eta^4}{V(U_+^2+\varepsilon^2)^2}.
\]
Finally \(\norm{\Delta u}_2^2\le B_+^2\), so this implies
\[
\mathcal Q_{\varepsilon,\alpha}[u]
\ge\frac{\eta^4}{VB_+^2(U_+^2+\varepsilon^2)^2}
\norm{\Delta u}_2^2.
\]
The smaller of the two coefficients proves
\eqref{eq:explicit-c}. If \(\Delta u=0\), the assertion is
immediate in either case.
\end{proof}

\begin{remark}[Asymptotic meaning]
This effective constant is generally conservative. It can deteriorate
as \(V\) grows, \(\lambda_1\to0\), \(H_\star\) grows, or the device
geometry changes. Thus the theorem does \emph{not} establish an
input-length-uniform physical resource bound. Its merit is that every
dependence is now stated as an explicit spectral sum or model parameter.
\end{remark}

\section{A detector corollary for the actual Hamiltonian evolution}
Let \(u(t)\) be a sufficiently smooth periodic solution of
\eqref{eq:pde} on \(0\leq t\leq T\), conserving positive mass \(m\)
and Hamiltonian \(H_0\). Let \(\chi\in H^4(\Omega;\mathbb C)\) have
zero spatial mean and write
\[
Z_\chi(t)=\int_\Omega\overline\chi\,u(t)\dd,\qquad
\delta=|Z_\chi(T)-Z_\chi(0)|.
\]
Let \(\lambda_1>0\) be the smallest positive eigenvalue of
\(-\Delta\), and use the amplitude bound \(U\) from the conserved
resources. For
\[
K_\chi=
\frac{\norm{(\Delta+\kappa)^2\chi}_2+
(3\beta U^2+5\gamma U^4)\norm{\chi}_2}{\lambda_1},
\]
the periodic equation, the polynomial nonlinear Lipschitz estimate,
and the mean-zero Poincar\'e inequality imply
\begin{equation}
\label{eq:obs}
\delta^2\leq K_\chi^2\,T
\int_0^T\norm{\Delta u(t)}_2^2\,dt.
\end{equation}
\begin{corollary}[Node-tolerant finite-time density-curvature cost]
\label{cor:transport}
If \(T>0\), \(K_\chi>0\), \(\delta>0\), and the smooth solution obeys
the conserved mass and Hamiltonian hypotheses above, then
\begin{equation}
\boxed{\quad
\int_0^T\mathcal Q_{\varepsilon,\alpha}[u(t)]\,dt
\geq\frac{c_\star\,\delta^2}{K_\chi^2\,T}.\quad}
\end{equation}
This remains valid if \(u(t,x)=0\) at some points and times.
\end{corollary}
\begin{proof}
The conserved quantities place every \(u(t)\) in the \emph{same}
class \(\mathcal A_{m,H_0}\), so Theorem~\ref{thm:coercivity}
gives one uniform \(c_\star>0\) for all \(t\). Integrate
\eqref{eq:coercive} in time and apply \eqref{eq:obs}.
\end{proof}

\section{Discrete Lean certificate and exact limitations}
On a two-node connected graph with real field values \(a,b\), set
\(m=a^2+b^2>0\), and use the simplified density weight
\(w_e(s)=s/(s+e)^2\), \(e>0\).
The two graph-Laplacian values are \(b-a,a-b\). Since
\(w_e(a^2)+w_e(b^2)\geq m/(m+e)^2\),
\begin{equation}
\boxed{\quad
(a-b)^2[w_e(a^2)+w_e(b^2)]
\geq\frac{m}{2(m+e)^2}\,[2(a-b)^2].\quad}
\end{equation}
The estimate holds even if \(a=0\) or \(b=0\), and the
coefficient is strictly positive when \(m,e>0\).
The separate Lean module
\texttt{PVNP/FixedMassNodal.lean} verifies this exact finite result
and algebraic detector-to-cost transfer. It does \emph{not}
kernel-formalize the Sobolev compactness proof of
Theorem~\ref{thm:coercivity} nor the continuum PDE integration and
conservation arguments in Corollary~\ref{cor:transport}.

\paragraph{Computational limitation.}
Neither \(c_\star\) nor \(K_\chi\) is shown to be independent of a
growing device domain, input length, resolution, or detector family.
The estimates concern a restricted observable of one declared PDE,
not all admissible SAT solvers. In particular, no unconditional
\(P_{\mathrm{WCC}}\ne NP_{\mathrm{WCC}}\) or classical \(P\ne NP\)
claim follows.

\paragraph{Source provenance.}
This supplemental note uses the maintained
\href{https://github.com/rickyjreyes/wct-pde/blob/main/src/wct_pde/projection_free.py}
{projection-free Hamiltonian equation} and follows the independent
\href{https://github.com/rickyjreyes/PVNP/blob/main/pvnp/proofs/phase-sensitive-observability.md}
{phase-observability study}. The original PVNP trilogy and its
published versions are not modified here.
\end{document}
